\section{Introduction}\label{sec:introduction}

How deeply must Kleene stars be nested to describe a regular language when
complement is also available? The answer depends on the permitted
operations. In this paper an expression over a fixed finite alphabet
$\Sigma$ starts from $0$, $1$, and the letters of $\Sigma$, denoting
$\varnothing$, $\{\eps\}$, and the corresponding singleton languages.
It uses union, concatenation, complement in $\Sigma^*$, and Kleene star.
Its height is given by
\[
\begin{aligned}
 \sht(0)=\sht(1)=\sht(a)&=0,\\
 \sht(P\cup Q)=\sht(PQ)&=\max\{\sht(P),\sht(Q)\},\\
 \sht(\neg P)&=\sht(P),&
 \sht(P^*)&=1+\sht(P).
\end{aligned}
\]
The \emph{generalized star height} $h_\Sigma(L)$ is the minimum height of
an expression defining $L\subseteq\Sigma^*$. Finite Boolean operations
preserve any common height bound. Finite languages have height zero, as
does $\Sigma^*=\neg0$. The complement operation keeps the height of its
operand; it does not remove stars already nested in that operand.

\begin{theorem}\label{thm:main}
For every finite alphabet $\Sigma$ and every regular language
$L\subseteq\Sigma^*$, one has $h_\Sigma(L)\le4$.
\end{theorem}

We prove the theorem by a complete construction over the given alphabet.
The bound is independent of $\Sigma$, $L$, and the number of states in a
recognizing automaton. Our finite choices can be found by exhaustive
search; we assert no useful bound on their sizes or on expression length.

\subsection{The problem and the three constructions}

The ordinary problem arose from expressions without complement. Eggan
related their star height to transition graphs~\cite{Eggan1963};
Dejean and Sch\"utzenberger established unbounded ordinary star height over
a binary alphabet~\cite{DejeanSchutzenberger1966}. Those lower bounds do
not carry over when complement is added. In the Boolean model,
Sch\"utzenberger characterized the height-zero, or star-free, languages
as exactly those recognized by aperiodic finite
monoids~\cite{Schutzenberger1965}.

Height-one results extend well beyond the aperiodic case. Henneman proved
that every language recognized by a finite commutative group has height
at most one, as recalled in \cite[Theorem~3.4]{PST1992}. Pin, Straubing,
and Th\'erien extended this conclusion to finite nilpotent groups of class
two~\cite[Theorem~7.3]{PST1992}. These algebraic results do not supply a
uniform bound for arbitrary finite monoids. Absolute boundedness and the
assertion that height one always suffices are different questions;
Straubing explicitly formulates both~\cite[p.~537]{Straubing2002}.
Theorem~\ref{thm:main} supplies an absolute bound. It does not exhibit a
language of height greater than one or decide whether one always suffices.

The alphabet convention is part of this problem. An alphabetic morphism
sends each letter to a letter or the empty word. Pin, Straubing, and
Th\'erien prove that inverse alphabetic morphisms do not increase generalized
height~\cite[Corollary~4.7]{PST1992}. They also express every regular
language as the inverse image, under a general morphism, of a language
of ordinary height at most one~\cite[Theorem~8.1]{PST1992}. Preservation of every
generalized-height bound under arbitrary inverse morphisms would therefore
already imply the bound one. Our component expressions are instead constructed
directly over the given alphabet; finite tags index tests and never become
new input letters.

This paper and two companions give different ways to bound the nesting.
\emph{Finite Monoid Computations and a Uniform Generalized Star-Height
Bound}~\cite[Theorem~1.1]{Uniform13} gives height thirteen using prediction
trials and a computation that records the full history of earlier tests.
The present construction gives height four using an affine correction and
several scales at which product boundaries may move. The companion
\emph{Generalized Star Height at Most Three}~\cite[Theorem~1.1]{Height3}
gives the stronger numerical bound three through arbitrary-function shift
tables and a decoder that controls its entire set of candidate cuts.
Each paper proves its own required algebra, timing statements, and
original-alphabet compilation. The purpose here is to develop the complete
multiscale method, including its two different residual decoders.

The split and affine-recovery framework is also developed in
\cite[Sections~2--4]{Uniform13}. Our local split lemma permits both factors
to have height at most one, and we prove that precise form below. The
finite probability argument behind our affine correction is an application
of the dependency-graph method of Erd\H{o}s and Lov\'asz
\cite[Section~2]{ErdosLovasz1975}. Our finite-window markers use the
marker-versus-periodicity principle associated with Krieger's marker lemma
\cite[Lemma~2]{Krieger1982}; the particular local rule and every periodic
transport claim needed here are proved in full. These comparisons concern
methods and mathematical scope: neither companion's numerical theorem is
a premise of this proof.

\subsection{From products to three splits}

A deterministic automaton turns each word into a transformation of a finite
set. These transformations form a finite monoid $M$, and their product in
reading order determines acceptance. We therefore seek low-height tests
for the value of a morphism $T:\Sigma^*\to M$. Cancellation and
commutativity are not assumed.

The word will be parsed into nonempty pieces called \emph{episodes},
followed by a remainder with no episode prefix. Each episode has a unique
end, determined by a bounded marker search or by the first departure from
a periodic continuation. A graph test records the input and output of a
fixed finite-state update on a sequence of episodes. To construct such a
test, we split selected episodes into a prefix test and a suffix test,
and use an already available graph test to skip the other episodes. The
two factors choose their finite information independently. Soundness must
therefore hold for every accepted pair, including pairs with different
guesses and a proposed suffix that ends too soon or too late.

The first difficulty is the product of the interval containing the cut:
neither side can read it in full. Suppose a segment has interval products
$(d_1,\ldots,d_l)\in M^l$ and carries a state in a finite vector space.
We choose one correction function $\beta$ on these tuples so that, for
every input state and every tuple, the corrected output is constant when
one coordinate $d_j$ varies. The prefix can then transmit the input state
and the earlier products; the suffix tests constancy while supplying the
later products. The unknown product $d_j$ need not be transmitted. Crucially,
the same $\beta$ works simultaneously for all inputs and tuples.

The second difficulty is timing. A suffix knows where it begins but may
not know the origin of the episode that produced it. A finite clock
restricts the possible pairs of prefix and suffix information to matching
tags or one uniquely determined mismatch. A deterministic affine
correction handles that mismatch. For matching tags, we choose boundaries
that remain fixed near markers and can move through a periodic region
without changing ordered monoid products. Several scales of movement and
a polynomial-indexed set of cut positions make the number of ambiguous
origin pairs smaller than the number of copies of every needed cut.

There are two remaining timing obstructions: failure of the clock or of
all inner scales, and instability of the later end search. Two buffers
handle them by different decoders. One first forces a common exact end and
then decodes an origin; the other first decodes an origin using bounded
data and then checks that origin's exact end. Their physical order lets
each suffix read the complete later update table. Those tables are defined
at every hypothetical monoid value, including impossible prefix values.
Lifting them to linear maps on a larger basis permits recovery on whole
episode sequences.

These operations lead to three applications of one split identity. Starting
from the empty skipped class, their heights are $0\to2\to3\to4$; the
intervening affine and basis recoveries use only finite Boolean operations.
To justify this count we must also control every component language on its
own, under false guesses as well as true ones. Each has only one unbounded
periodic continuation and hence a finite cover by word templates
$a b^h c$. Ultimately periodic tests on $h$ compile directly over $\Sigma$.
No tag or state is introduced as an extra input letter.

\paragraph{Organization.}
Section~\ref{sec:algebra} gives the monoid reduction, independent-witness
split principle, and affine-line lemma. Section~\ref{sec:clock} constructs
the finite clock and proves its roster-independent bound.
Section~\ref{sec:scales} builds the roster, scales, and moving boundaries
and counts ambiguous origins. Section~\ref{sec:episodes} places all windows
and proves exact end detection. Section~\ref{sec:splits} defines the total
updates and performs the three splits. Section~\ref{sec:compilation}
compiles the individual tests, handles the remainder, and completes the
height calculation.
